\documentclass{amsart}
% For dealing with references we use the comment environment
\usepackage{verbatim}
\newenvironment{reference}{\comment}{\endcomment}
%\newenvironment{reference}{}{}
\newenvironment{slogan}{\comment}{\endcomment}
\newenvironment{history}{\comment}{\endcomment}

% For commutative diagrams we use Xy-pic
\usepackage[all]{xy}

% We use 2cell for 2-commutative diagrams.
\xyoption{2cell}
\UseAllTwocells

% We use multicol for the list of chapters between chapters
\usepackage{multicol}

% This is generally recommended for better output
\usepackage{lmodern}
\usepackage[T1]{fontenc}

% For cross-file-references

% Package for hypertext links:
\usepackage{hyperref}

% For any local file, say "hello.tex" you want to link to please

% Theorem environments.
%
\theoremstyle{plain}
\newtheorem{theorem}[subsection]{Theorem}
\newtheorem{proposition}[subsection]{Proposition}
\newtheorem{lemma}[subsection]{Lemma}

\theoremstyle{definition}
\newtheorem{definition}[subsection]{Definition}
\newtheorem{example}[subsection]{Example}
\newtheorem{exercise}[subsection]{Exercise}
\newtheorem{situation}[subsection]{Situation}

\theoremstyle{remark}
\newtheorem{remark}[subsection]{Remark}
\newtheorem{remarks}[subsection]{Remarks}

\numberwithin{equation}{subsection}

% Macros
%
\def\lim{\mathop{\mathrm{lim}}\nolimits}
\def\colim{\mathop{\mathrm{colim}}\nolimits}
\def\Spec{\mathop{\mathrm{Spec}}}
\def\Hom{\mathop{\mathrm{Hom}}\nolimits}
\def\Ext{\mathop{\mathrm{Ext}}\nolimits}
\def\SheafHom{\mathop{\mathcal{H}\!\mathit{om}}\nolimits}
\def\SheafExt{\mathop{\mathcal{E}\!\mathit{xt}}\nolimits}
\def\Sch{\mathit{Sch}}
\def\Mor{\mathop{\mathrm{Mor}}\nolimits}
\def\Ob{\mathop{\mathrm{Ob}}\nolimits}
\def\Sh{\mathop{\mathit{Sh}}\nolimits}
\def\NL{\mathop{N\!L}\nolimits}
\def\CH{\mathop{\mathrm{CH}}\nolimits}
\def\proetale{{pro\text{-}\acute{e}tale}}
\def\etale{{\acute{e}tale}}
\def\QCoh{\mathit{QCoh}}
\def\Ker{\mathop{\mathrm{Ker}}}
\def\Im{\mathop{\mathrm{Im}}}
\def\Coker{\mathop{\mathrm{Coker}}}
\def\Coim{\mathop{\mathrm{Coim}}}

% Boxtimes
%
\DeclareMathSymbol{\boxtimes}{\mathbin}{AMSa}{"02}

%
% Macros for moduli stacks/spaces
%
\def\QCohstack{\mathcal{QC}\!\mathit{oh}}
\def\Cohstack{\mathcal{C}\!\mathit{oh}}
\def\Spacesstack{\mathcal{S}\!\mathit{paces}}
\def\Quotfunctor{\mathrm{Quot}}
\def\Hilbfunctor{\mathrm{Hilb}}
\def\Curvesstack{\mathcal{C}\!\mathit{urves}}
\def\Polarizedstack{\mathcal{P}\!\mathit{olarized}}
\def\Complexesstack{\mathcal{C}\!\mathit{omplexes}}
% \Pic is the operator that assigns to X its picard group, usage \Pic(X)
% \Picardstack_{X/B} denotes the Picard stack of X over B
% \Picardfunctor_{X/B} denotes the Picard functor of X over B
\def\Pic{\mathop{\mathrm{Pic}}\nolimits}
\def\Picardstack{\mathcal{P}\!\mathit{ic}}
\def\Picardfunctor{\mathrm{Pic}}
\def\Deformationcategory{\mathcal{D}\!\mathit{ef}}

\setlength{\emergencystretch}{3em}
\begin{document}
\title[Stacks Project, Tag 0CQY]{412 --- Derived complete power-torsion modules have bounded torsion}
\maketitle
\noindent Copyright (C) 2005--2025 Johan de Jong. Stacks Project authors. Source: \href{https://stacks.math.columbia.edu/tag/0CQY}{Tag 0CQY}.

\noindent Editorial difficulty note (not author text). Difficulty H3: The proof reduces to a nonzerodivisor principal ideal and represents the module by a quotient of complete torsion-free modules. A Baire-category argument and an open-mapping theorem promote pointwise power torsion to one uniform annihilating power..

\noindent Editorial provenance note (not author text). History: excerpt from revision \texttt{a0444\allowbreak 6e57e\allowbreak c1fbc\allowbreak 252a8\allowbreak 71afc\allowbreak ec775\allowbreak 2fb28\allowbreak 07b14}, more-algebra.tex, lines 27014--27064. Reference presentation was adapted; original-author context and core support blocks were retained or added. The mathematical wording is unchanged. Licensed under GNU FDL 1.2 or later; no invariant sections or cover texts. See ../licenses/COPYING. Context blocks reproduce author definitions and supporting results; they are not additional counted problems.

\begin{lemma}[Bhatt]
\label{lemma-torsion-and-derived-complete}
Let $I$ be a finitely generated ideal in a ring $A$.
Let $M$ be a derived complete $A$-module. If $M$ is
an $I$-power torsion module, then $I^nM = 0$ for some $n$.
\end{lemma}

\begin{proof}
Say $I = (f_1, \ldots, f_r)$. It suffices to show
that for each $i$ there is an $n_i$ such that $f_i^{n_i}M = 0$.
Hence we may assume that $I = (f)$ is a principal ideal.
Let $B = \mathbf{Z}[x] \to A$ be the ring map sending $x$ to $f$.
By Lemma \href{https://stacks.math.columbia.edu/tag/0924}{lemma restriction derived complete (Tag 0924)}
we see that $M$ is derived complete as a $B$-module
with respect to the ideal $(x)$. After replacing $A$ by
$B$, we may assume that $f$ is a nonzerodivisor in $A$.

\medskip\noindent
Assume $I = (f)$ with $f \in A$ a nonzerodivisor.
According to Example \href{https://stacks.math.columbia.edu/tag/09AT}{example derived complete modules (Tag 09AT)}
there exists a short exact sequence
$$
0 \to K \xrightarrow{u} L \to M \to 0
$$
where $K$ and $L$ are $I$-adically complete $A$-modules
whose $f$-torsion is zero\footnote{For the proof it is enough
to show that there exists a sequence $K \xrightarrow{u} L \to M \to 0$
where $K$ and $L$ are $I$-adically complete $A$-modules. This can
be shown by choosing a presentation $F_1 \to F_0 \to M \to 0$
with $F_i$ free and then setting $K$ and $L$ equal to the
$f$-adic completions of $F_1$ and $F_0$. Namely, as $f$
is a nonzerodivisor these completions will be the
derived completions and the sequence will remain exact.}.
Consider $K$ and $L$ as
topological modules with the $I$-adic topology. Then $u$ is continuous.
Let
$$
L_n = \{x \in L \mid f^n x \in u(K)\}
$$
Since $M$ is $f$-power torsion we see that
$L = \bigcup L_n$. Let $N_n$ be the closure of $L_n$ in $L$.
By Lemma \href{https://stacks.math.columbia.edu/tag/0CQV}{lemma consequence baire complete module (Tag 0CQV)}
we see that $N_n$ is open in $L$ for some $n$. Fix such an $n$.
Since $f^{n + m} : L \to L$ is a continuous open map, and since
$f^{n + m} L_n \subset u(f^m K)$ we conclude that
the closure of $u(f^mK)$ is open for all $m \geq 1$.
Thus by Lemma \href{https://stacks.math.columbia.edu/tag/0CQW}{lemma open mapping (Tag 0CQW)}
we conclude that $u$ is open. Hence $f^tL \subset \Im(u)$
for some $t$ and we conclude that $f^t$ annihilates $M$
as desired.
\end{proof}
\end{document}
